% Nouveau EGA V (cote 25): archive cote 25, site page 75 (PDF page 76).
% typescript. Draft compared: none.
% AI-transcribed by Claude Opus 5.5 (claude-opus-5-5, adaptive thinking, effort high), Message Batches API, batch ega25b (msgbatch_01A7n5YuXubUPNK1zsepi8RG), 2026-11-13; not yet checked by eye.
\scanpage{25/75}
\mspage{70}

20.13. \uline{Morphismes élémentaires et théorème de M. Artin}.

\uline{Définition} 13.1. Un morphisme $f:X\to Y$ de préschémas est dit ``morphisme
\xout{\unsure{élém}} élémentaire'' si X est \xout{\unsure{un}} Y-isomorphe à un préschéma de la forme
$X' - Z$ , où $X'\to X$ est un morphisme lisse, projectif, à fibres géomé-%
triques connexes\xout{\unsure{,}} de dimension 1, et où Z est un sous-préschéma fermé
de X', \xout{\unsure{étale sur}} tel que le morphisme $Z\to Y$ soit étale, surjectif, et
de degré constant. Un morphisme est dit \xout{\unsure{multip-élémentaire}} poly-élémen-%
taire s'il est un composé de morphismes élémentaires. Un préschéma X
sur un corps k est dit polyélementaire (sur k), si le morphisme structu-%
ral $X\to \mathrm{Spec}(k)$ est poly-élémentaire.

\uline{Théorème} 13.2. \inserted[typed]{(M.Artin)} Soient X un préschéma \inserted[typed]{géom. irréd.} sur un corps k \inserted[typed]{parfait infini}, x un point lisse
de X, alors x admet un système fondamental de voisinages ouverts
polyélémentaires.

Quitte à remplacer X par un voisinage donné de x, \deleted{on \xout{\unsure{peutxsuppo}} es} \inserted{il revient au}
\deleted{ramené à} \inserted{même de} montrer qu'il existe un voisinage ouvert élémentaire de x dans X.
\deleted{On peut supposer évidemment X lisse, et} \inserted[typed]{\deleted{et affine}} \over{r}{R}aisonnant par récurrence sur
la dimension n de X, on \xout{\unsure{peutxsupposer}} est ramené à montrer \inserted[typed]{que si $n>0$ ,} \deleted{qu'}il existe
un voisinage ouvert U de x et un morphisme élémentaire $f: U\to V$,
V étant un schéma lisse sur k (nécessairement géom. irréd \inserted[typed]{et} de dim n-1).
\xout{\unsure{Lorsque n=1}} (Le cas n=0 est évidemment trivial, étant donné qu\over{e}{'} alors
X est k-isomorphe à Spec(k), qui est polyélémentaire sur Spec(k), étant
entendu dans 13.1. qu'on n'exc\hand{l}ut pas le composé de la famille vide de
morphismes - on aurait dû le signaler sous une forme ou une autre dans
13.1. ). La nécessité de supposer k parfait au préalable apparait déjà
dans le cas n=1: alors on prend pour X' le modèle projectif normal
canonique du corps des fonctions K de X (cf Chap II) par.7), le fait que

\mspage{75}
